\subsection{User-Requested Nearest-Time and Final-Occupant Scenario}\label{sec:tail-policy}

Engine 1.4.0 adds a separately configured behavioural scenario. It does not replace the historical 1.3.0 LED/barrier experiment below. In the current emergency preset, CM1--CM4 retain their posts but broadcast ``use the nearest exit by estimated travel time'' instead of a fixed portico destination. CM5 and CM6 retain their balcony instructions. Each compliant listener evaluates known, feasible exits from the next reachable junction using Eq.~\eqref{eq:perceived}, with current visible queues and individual speed. The unfinished link remains mandatory. This common public rule does not supply unknown exits or remote queue information. Ordinary current-state switching inertia still applies.

The optional \texttt{tailFallback} rule begins when at most 40 people remain anywhere inside the building. Exterior wheelchair waiting is not included in this count. A person who has completed a link and remained in its discharge queue for strictly more than 10 continuous seconds during this final phase selects
\[
g_i^*=\arg\min_{g\in K_i\cap\mathcal F_i(t)}\widehat T_{ig}(t),
\]
where $K_i$ is the known-exit set and $\mathcal F_i(t)$ restricts routes by accessibility, width, learned closures and the existing return-lane convention. The snapshot uses current waiting occupants and the original visibility rule. Foyer preference, fixed random taste, companion destinations, prior guidance and the eight-second switching threshold do not enter this final choice. Ties are broken by exit identifier.

The timer starts at the later of entry into the final phase and the current queue arrival. Progress onto a new link resets it. Recognition, preparation and seat/group waiting are not locomotion stalls. A person already traversing an exterior-crossing link continues outside. Each stalled person acts individually; moving members of the final forty are not redirected by someone else's stall.

The selected exit and complete route are then locked. Subsequent guidance, companion choices and changing queue estimates are ignored, while travel length, FIFO, lane direction, gate service and finite storage remain unchanged. A final-commitment event retains the decision state, alternatives and route. If no feasible known exit exists, none is invented. A locked path encountering an obstruction can remain unresolved; the lock cannot authorize crossing a barrier or manufacture capacity.

These thresholds and permanent withdrawal from guidance are user-requested scenario assumptions, not values calibrated to Bailey or established by the literature. The emergency JSON enables them explicitly, and the interface allows them to be enabled for other scenarios. Omission disables the rule, preserving historical configurations. New regression experiments are reported separately; the historical 100-replication-per-policy means below must not be attributed to this revision.


A regression experiment evaluated ten paired audiences (seeds 5001--5010) with the revised manager messages, with and without the final-occupant rule. Both arms cleared all 1,329 people in every run, averaged 182.5 s and met 180 s in 3/10 replications. All paired clearance times were identical despite 90 final commitments in the enabled arm. All 90 retained an exit already minimizing their estimated time. This sample does not demonstrate a clearance benefit. It was not used to retune parameters and does not replace the historical 100-replication study.
